\chapter{Validation, Limitations, and Next Steps}

\section{What Has Been Demonstrated}

The project has progressed beyond a model-only experiment. The completed chain
includes a repeatable egocentric recording workflow, a browser-based labeling
application with structured exports, a curated engine-model dataset and
generated stress variants, controlled IndustReal and engine-model evaluation, a
causal live inference service with fixed-lag step commitment, and replay, live,
and HoloLens presentation paths.

The strongest supplied engine-model offline result is 96.21\% accuracy,
97.33 edit score, and 97.51 F1@50. The final online result is 94.12\%,
84.45 edit score, and 89.50 F1@50. The deployed robustness checkpoint reaches
90.57\% on the scrambled-order stress condition and starts producing output
after approximately 0.2 seconds in the reported gate.

\section{Current Limitations}

\begin{enumerate}
    \item The current held-out split is recording-disjoint but not
    subject-disjoint; performance on new operators is not yet isolated.
    \item The model recognizes procedure-step evidence. It is not a complete
    visual inspection system and should not be interpreted as proof that a
    part is installed correctly.
    \item The live output has a deliberate delay from clip context and decoder
    stabilization.
    \item The wearable client currently presents a screen-space overlay rather
    than a spatially anchored instruction at the physical part.
    \item The learned vocabulary covers ten modeled parts, while a presentation
    or overlay may show additional workflow or verification states. These
    should not be confused with learned recognition classes.
    \item The feature encoders are computationally substantial. A production
    deployment would need a defined edge, workstation, or network-inference
    architecture and an operational monitoring plan.
\end{enumerate}

\section{Recommended Next Steps}

\begin{enumerate}
    \item Collect more operators, lighting conditions, camera placements, and
    procedure variations; introduce a subject-disjoint evaluation split.
    \item Expand labels beyond completion events to include wrong part, omitted
    step, rework, hesitation, and verification outcomes.
    \item Improve the live path by measuring per-stage latency, testing
    hardware-aware optimization, and defining an acceptable alert-delay
    envelope with end users.
    \item Add spatial registration and part localization only where the
    operator workflow benefits from an anchored visual cue.
    \item Establish a production evaluation protocol with confidence
    thresholds, false-alert review, model/version traceability, and approved
    human decision points.
\end{enumerate}

\section{Reproducibility and Dependencies}

The repository contains the model code, dataset preparation scripts, recorded
assets, checkpoint, and the replay/live services. The software path is based on
Python and PyTorch, with FastAPI/Uvicorn for the services and standard video
processing utilities for frame handling. CUDA is used when available; CPU
execution is supported for inspection but is not the intended performance
configuration for the full encoder path. The HoloLens presentation is a
separate Unity UWP project using MRTK/OpenXR. The reproducibility guide in the
project defines the staged sequence: prepare features and labels, generate
stress data, train, run the acceptance gate, and launch the served applications.
